\documentclass[11pt]{article}

% Axioms for physical reasoning: codifying the Seiberg-Witten solution in Lean.
% Companion working paper to mrdouglasny/seiberg-witten and the ICML 2026 talk.
% Self-contained; compiles with pdflatex (no exotic packages).

\usepackage[margin=1in]{geometry}
\usepackage{amsmath,amssymb}
\usepackage{microtype}
\usepackage[hidelinks]{hyperref}
\usepackage{xcolor}

\newcommand{\C}{\mathbb{C}}
\newcommand{\Z}{\mathbb{Z}}
\newcommand{\HH}{\mathbb{H}}
\newcommand{\Sp}{\operatorname{Sp}}
\newcommand{\Imag}{\operatorname{Im}}
\newcommand{\code}[1]{\texttt{#1}}

\title{Axioms for physical reasoning:\\ codifying the Seiberg--Witten solution in Lean}
\author{Michael R.\ Douglas\\[2pt]
  \small Center of Mathematical Sciences and Applications, Harvard University}
\date{\small Working paper, draft. Companion to the formal development at
  \href{https://github.com/mrdouglasny/seiberg-witten}{\code{mrdouglasny/seiberg-witten}}
  and to the ICML 2026 talk \emph{Validation of AI-generated results in theoretical physics}.}

\begin{document}
\maketitle

\begin{abstract}
Most results in theoretical physics are not theorems. They are conclusions of arguments the
community trusts without a proof from first principles. For the deepest of them a rigorous
proof is conceivable but lies far out of reach---it would presuppose a nonperturbative
construction of the interacting quantum field theory, a goal at least as hard as the
Yang--Mills mass-gap Millennium problem. We argue that an interactive theorem prover is the
right instrument for such a result even so: not to \emph{prove} it---that is beyond
reach---but to state the physical reasoning as a short list of \textbf{explicit, named
axioms}, prove every mathematical consequence of them, and keep a machine-checked line between
the two. The trust that remains then rests on that short, inspectable list, and the prover can
report, for any downstream result, exactly which assumptions it used. We carry this out for
the Seiberg--Witten solution of $\mathcal{N}=2$ $SU(N)$ super-Yang--Mills, formalized in Lean
4. We describe what is proved, what is assumed, how the assumptions are checked---external
review and an independent numerical oracle---and why this discipline is a sound standard for
validating AI-generated results in theoretical physics.
\end{abstract}

\section{Physics is mostly not theorems}

A mathematician who formalizes a theorem has a proof to formalize. A physicist who formalizes
a \emph{result} usually does not. The Seiberg--Witten (SW) solution is the textbook case. No
rigorous proof says that the low-energy physics on the Coulomb branch of $\mathcal{N}=2$
$SU(N)$ gauge theory is governed by the curve $y^2 = P_N(x)^2 - \Lambda^{2N}$. One is
\emph{conceivable}---but it would first require constructing the interacting $\mathcal{N}=2$
theory nonperturbatively, a problem at least as hard as the Yang--Mills mass-gap Millennium
problem, and so out of reach for the foreseeable future. What we have instead is an
\emph{argument}: a small set of physical assumptions---holomorphy from supersymmetry, the
exact pattern of singularities, electric--magnetic duality, the weak-coupling
asymptotics---that together pick out that curve, believed because three independent
computations of the same quantity agree.

So ``formalizing Seiberg--Witten'' cannot mean ``proving it''---not today, and not without
first solving a Millennium-scale problem. That seems to put the result beyond formal methods.
We argue the reverse: such a result is precisely where formalization helps most, \emph{provided
one is honest about what it does}.

\begin{quote}
State the physical assumptions as \textbf{explicit, named axioms}; prove every
\textbf{mathematical consequence} of them; and keep a machine-checked line between the two.
\end{quote}

We call this second category---the non-rigorous reasoning physicists nonetheless
trust---\textbf{foundations for physical reasoning}, by analogy with foundations of
mathematics. A formalization then does three things no informal write-up can: it makes the
dependency structure machine-checkable, it concentrates the residual trust onto a few named
axioms, and it lets the kernel detect a contradiction among them.

This matters now because of \emph{who}, or what, increasingly produces such arguments. As
large language models begin to propose physical and mathematical derivations, their
characteristic failure is not a wrong algebra step a referee would catch but a
\emph{hallucinated assumption}---an unphysical premise, or a dropped topological subtlety,
presented as though derived. A machine-checked boundary between assumed and proved is precisely
what keeps such a system honest; that is why we frame the discipline below as a standard for
validating AI-generated physics, and return to it in \S6.

\section{Two kinds of formalization in physics}

\begin{center}
\renewcommand{\arraystretch}{1.3}
\begin{tabular}{@{}p{0.26\textwidth}p{0.30\textwidth}p{0.34\textwidth}@{}}
\hline
& \textbf{What it is} & \textbf{Examples} \\
\hline
\textbf{Mathematical physics} & A rigorous statement exists; you prove it. & Stability of
matter; constructive QFT; the lattice-gauge mass gap---\emph{and the proved spine of this
project}: the curve, its genus, special-K\"ahler positivity, the period matrix. \\
\textbf{Foundations for physical reasoning} & Codifying the \emph{non-rigorous} reasoning the
community trusts, where no proof is within reach (a rigorous one would be Millennium-hard). &
\emph{The Seiberg--Witten argument itself.} \\
\hline
\end{tabular}
\end{center}

\noindent This is not a value judgment but a bookkeeping discipline, and it is necessary
because the kernel does not enforce it. A theorem can be proved, axiom-free, and still be the
\textbf{wrong statement}---a dropped hypothesis, a flipped sign, a domain quietly narrowed. And
a physical assumption written as an innocent-looking definition can smuggle a falsehood past
the kernel entirely. The discipline forces every physical input into the open as a named
axiom, and every mathematical claim onto the proved side.

\section{The Seiberg--Witten solution in one paragraph}

\emph{For a reader outside physics: the SW solution maps the exact low-energy behaviour of a
quantum system to the geometry of a specific family of algebraic curves. What matters for the
formalization is less the physical meaning of the terms below than the rigid mathematical data
they impose.}

On the Coulomb branch---coordinates $u_2,\dots,u_N$, rank $r = N-1$---the low-energy theory is
abelian $U(1)^r$ gauge theory. Its data are the special coordinates $a_i$ and their duals
$a_{D,i}$ (periods of a meromorphic differential), the coupling matrix $\tau_{ij}$, and the
central charge $Z = n_e\cdot a + n_m\cdot a_D$ of a BPS state, whose modulus is its mass. The
SW claim is that all of this \emph{is} the period geometry of $y^2 = P_N(x)^2 - \Lambda^{2N}$
($P_N$ a degree-$N$ polynomial in a coordinate $x$, distinct from the moduli $u$). The
assumptions that single out that curve are: holomorphy / non-renormalization (one holomorphic
prepotential $\mathcal{F}$, one-loop-exact perturbation theory); the exact singularity
structure (the branch is singular exactly where a BPS state goes massless, and nowhere else);
electric--magnetic duality ($\Sp(2r,\Z)$ monodromy); the $\Z_{2N}$ $R$-symmetry fixing the
singularities; and matching to the one-loop-plus-instanton expansion at weak coupling.

\section{The formalization}

\subsection{What is proved}

On a hyperelliptic-curve / period / Jacobian library, the following are theorems---sorry-free,
and (except where noted) depending only on the three standard logical axioms:

\begin{itemize}
\item the $SU(N)$ curve has \textbf{genus $N-1$};
\item its period matrix $\tau$ is \textbf{symmetric with $\Imag\tau \succ 0$}---the positivity
of the special-K\"ahler metric, i.e.\ the physical no-ghost condition;
\item for $SU(2)$, the curve is singular at \textbf{exactly $u = \pm\Lambda^2$}---the monopole
and dyon points, and no others (this is a discriminant computation, not lore);
\item the Picard--Lefschetz monodromy of a vanishing cycle is \textbf{symplectic}---the reason
the monodromy group sits in $\Sp(2r,\Z)$;
\item for the rank-1 base, the modular function $\lambda = \theta_2^4/\theta_3^4$ (built from
Mathlib's Jacobi theta functions) is a proved \textbf{$\Gamma(2)$-Hauptmodul}: invariant under
both generators of $\Gamma(2)$, with the cusp symmetry $1-\lambda = \theta_4^4/\theta_3^4$.
\end{itemize}

(Declaration names and exact axiom footprints are in the repository; we keep them out of the
prose here.)

\subsection{What is assumed}

The five physical assumptions of \S3 are stated as named axioms in one module, each with its
source and a faithfulness note. These are the irreducible physics---there is nothing to prove
about them. (We use Lean's global \code{axiom} keyword deliberately, as a
\emph{dependency-tracking} device: it is exactly what \code{\#print axioms} reports, so every
assumption surfaces natively in the footprint of any result that uses it. Because a false
axiom would make the whole system inconsistent, each is treated as a debt to discharge, not a
convenience---the audit tracks them as such.)

The headline claims are different in kind. That the SW curve \emph{realizes} a theory
satisfying the assumptions (existence), and that the assumptions fix the solution
\textbf{uniquely up to $\Sp(2r,\Z)$ duality} (uniqueness), are not physical inputs but
\emph{mathematical} consequences: given the assumptions, each should be a \textbf{theorem}---
uniqueness, in particular, is a period-map rigidity statement. We state them as named axioms
today only as placeholders, because their proofs are not yet in place (they need the covering
and Riemann-mapping machinery of \S7). Discharging them is precisely that program, and it
\emph{purifies} the axiom list: it moves these mathematical claims onto the proved side,
leaving only the five genuinely-physical inputs as assumptions. Until then, printing the axioms
behind a result shows both---and makes visible exactly which placeholder still stands between
the physics and the theorem.

For $SU(2)$ we have carried this out. The rank-1 uniqueness statement is now a
\textbf{theorem}: printing its axioms returns only the three logical axioms, the physics
monodromy input, and two named \textbf{classical-math} axioms---the modular-$\lambda$ covering
$\HH/\Gamma(2)\cong\C\setminus\{0,1\}$ and the uniqueness of its developing maps---with
\textbf{no bespoke placeholder}. The two tiers are now explicit in the proof itself: the
physics axiom is irreducible, while the two math axioms are debt, discharged once the
(Mathlib-scale) complex analysis of \S7 lands. This is the purification in miniature---the
headline reduced to \emph{physics $+$ standard mathematics}, with the machine confirming
nothing else slipped in.

\subsection{Two things formalizing forced into the open}

\paragraph{The type system can refuse a convenient falsehood.} Our first attempt made the
periods $a, a_D$ and the coupling $\tau$ single-valued functions of the moduli. External review
flagged this as \emph{vacuous}: a global holomorphic map from the moduli space to the Siegel
upper half-space is constant by Liouville's theorem---a free theory. The data must instead live
on the universal cover of the moduli space minus its singular locus. The corrected types encode
this, refusing the free-theory shortcut that prose slides over.

\paragraph{A surprise comes out derived, not assumed.} Argyres--Douglas points---where the
theory becomes an interacting CFT with no Lagrangian---look as though they need their own
postulate. In the formalization they are instead a \emph{definition}: the locus where two
\textbf{mutually non-local} BPS charges (nonzero Dirac pairing) go massless together, built
from the central charge and the charge lattice already in hand. (This locus is \emph{empty} for
pure $SU(2)$---the monopole and dyon go massless at the two different points
$\pm\Lambda^2$---but the definition is general, correctly accommodating the matter flavours and
higher-rank theories where AD points actually occur.) A phenomenon the theory should explain
falls out as a consequence---a check the axioms pass rather than a cost they incur.

\section{Checking the axioms: validation by independent agreement}

Concentrating trust onto named axioms is worthwhile only if those axioms are then checked. We
use three tools, cheapest first.

\paragraph{An independent numerical oracle.} Where a statement has computable content, we
evaluate it in a separate engine (mpmath / numpy) that shares no code or authorship with the
Lean. A run of 44 checks at 40-digit precision confirms the assumed theta-function identities
($\theta_3^4 = \theta_2^4 + \theta_4^4$, $\theta_3 \neq 0$), verifies that the proved $\lambda$
transformation laws are \emph{faithful} (catching any sign or shift error), anchors the
normalization ($\lambda(i) = \tfrac12$), and checks the $SU(2)$ singularity locus. Tellingly,
the oracle first caught a bug \textbf{in itself}---a fractional power evaluated on the wrong
branch---before agreeing with the Lean to full precision: the method earning its keep. A
numerical pass is necessary, not sufficient; it validates values, not logical structure (it
cannot see a too-strong hypothesis or a vacuous one).

\paragraph{Cross-model and self-review} against the pinned sources, recorded per result,
separating ``is it proved?'' from ``does it say what the source says?''.

\paragraph{The physicist's triangle.} The deepest reason the SW lore is trusted is that three
independent routes to the prepotential agree: weak-coupling perturbation theory, the instanton
expansion (Nekrasov's partition function, itself later made rigorous), and the exact periods of
the curve. This independent agreement is the empirical content behind the axioms, and our
protocol attaches it as a cross-check rather than asserting it.

\section{Why this is the right standard for AI-generated physics}

When an AI produces a physics result, the dangerous failure is rarely a wrong algebra step a
referee would catch. It is an \textbf{unstated assumption}---a \emph{hallucination}, in
machine-learning terms---slipped in as though derived. Reasoning from named physical axioms
turns the kernel and type system into a guardrail that makes this impossible to hide:

\begin{itemize}
\item a physical input must be \emph{named} to be used;
\item every named input shows up when the prover lists the axioms behind a result;
\item an honest definition---the universal-cover carrier, Argyres--Douglas as a derived
locus---cannot encode a convenient falsehood, because the kernel checks the types;
\item and \emph{justifying} the axioms remains the physicist's job, done by the same
independent-agreement checks that earn trust in the first place.
\end{itemize}

Formalization does not make the physics true. It makes the reasoning \textbf{auditable}---which
is exactly what a result no one can prove from first principles needs.

\section{What remains, and where it lives}

The discipline also makes the remaining \emph{mathematical} debt legible. For $SU(2)$ the
uniqueness headline is, as just described, already a theorem modulo named classical-math
axioms; the debt is to \emph{discharge those axioms}, and it reduces to a few
sharply-identified, Mathlib-scale developments: the Jacobi triple product (which discharges the
theta identities and non-vanishing); the modular-$\lambda$ covering
$\HH/\Gamma(2)\cong\C\setminus\{0,1\}$ (whose invariance half is already proved, leaving the
surjectivity); and the Gauss--Manin / variation-of-Hodge-structure layer over the moduli space.
The covering rests in turn on the \textbf{Riemann mapping theorem} and the \textbf{Schwarz
reflection principle}---general complex analysis, absent from Lean (RMT is formalized in
Isabelle/HOL), which we have scoped as upstream contributions. The point is that ``is this
physics right?'' becomes a precise, prioritized list of theorems to build.

\section{Related work}

The assurance artifacts (correspondence index, fidelity review, numerical validation, axiom
audit) follow the conventions of the author's own \emph{formalization-assurance} project
(\code{math-commons/\discretionary{}{}{}formalization-assurance}), separate from and predating this work; the
upstream mathematical contributions are framed for the Lean FRO's Tau Ceti library. The Riemann mapping theorem is formalized in
Isabelle/HOL (Paulson); the Schwarz reflection principle and the SW solution appear
unformalized before this work.

\section{Conclusion}

A rigorous proof of the Seiberg--Witten solution lies beyond reach---it presupposes
constructing the quantum field theory itself---yet the result is as certain as anything in
physics. A proof assistant can still treat it rigorously: by separating the machine-checked
consequences from the physical lore, naming the lore as auditable axioms, and validating those
axioms by independent agreement. Proved and assumed as separate, machine-checked columns; trust
narrowed to a short vetted list; faithfulness enforced by types and by an independent oracle:
this, we argue, is the right standard for validating AI-generated results in theoretical
physics.

\appendix

\section{The physics axioms (irreducible inputs)}

The physical content is carried by named axioms in \code{SeibergWitten/Physics/}. They are
currently \emph{uninterpreted predicates} over the carrier \code{LowEnergyData} on the
universal cover \code{SmoothCover} of $\C^r\setminus\Delta$ (the cover, not $\C^r$---the
Liouville point of \S4.3); they stand in for the full definitions and are the genuine physics,
not provable but only justified.

\begin{center}\small
\renewcommand{\arraystretch}{1.35}
\begin{tabular}{@{}p{0.34\textwidth}p{0.60\textwidth}@{}}
\hline
\textbf{Axiom} & \textbf{Physical hypothesis} \\
\hline
\code{SpecialGeometry} & \textbf{H1} Holomorphy / special geometry: $a_D=\partial\mathcal{F}/\partial a$, $\tau=\partial^2\mathcal{F}/\partial a^2$, one-loop-exact prepotential. \\
\code{BecomesMasslessAt} & \textbf{H2} A BPS state becomes massless ($Z\to0$) on the singular locus $\Delta$---the building block of ``singular $\Leftrightarrow$ massless''. \\
\code{PicardLefschetzAt\-Generic\-Singularities} & \textbf{H3} Monodromy around a generic point of $\Delta$ is the Picard--Lefschetz class $v\mapsto v+\langle v,n_a\rangle\,n_a$. \\
\code{EMDualityConsistent} & \textbf{H4} Electric--magnetic duality: $\Sp(2r,\Z)$ monodromy. \\
\code{RSymmetryConsistent} & \textbf{H5} The $U(1)_R\to\Z_{2N}$ $R$-symmetry fixes the positions of $\Delta$. \\
\code{AsymptoticMatch} & \textbf{H6} One-loop $+$ instanton expansion matches the exact periods at weak coupling. \\
\code{SameSWMonodromy} & ($SU(2)$) Two period maps carry the same $\Sp(2,\Z)$ monodromy---the physics input to the rank-1 rigidity theorem. \\
\hline
\end{tabular}
\end{center}

\noindent Structural (carrier / interface, not assertions): \code{SmoothCover},
\code{SmoothCover.proj} (the universal-cover carrier) and \code{DualityEquiv} (the
$\Sp(2r,\Z)$ frame-change equivalence).

\section{Mathematical axioms, used provisionally (debt)}

These are \emph{mathematical} claims, distinct in kind from the physics of Appendix A. Each is
true standard mathematics but not yet proved in our development; each is a debt, tracked
natively by \code{\#print axioms}.

\medskip
\noindent\textbf{B.1 --- Headline implications, provisionally axiomatized pending proof.}
Mathematical consequences of the physics axioms that \emph{should be theorems} (\S4.2):

\begin{center}\small
\renewcommand{\arraystretch}{1.35}
\begin{tabular}{@{}p{0.40\textwidth}p{0.54\textwidth}@{}}
\hline
\textbf{Axiom} & \textbf{Claim} \\
\hline
\code{sw\_effective\_theory\_\-unique\_up\_to\_duality} & The hypotheses fix the solution uniquely up to $\Sp(2r,\Z)$ duality. \textbf{Proved for $SU(2)$} (\code{sw\_su2\_unique}); general rank still provisional. \\
\code{sw\_curve\_admits\_\-effective\_theory} & The SW curve realizes a theory satisfying the hypotheses (existence). \\
\hline
\end{tabular}
\end{center}

\medskip
\noindent\textbf{B.2 --- Classical mathematics not yet in Mathlib, assumed pending upstream
formalization:}

\begin{center}\small
\renewcommand{\arraystretch}{1.35}
\begin{tabular}{@{}p{0.27\textwidth}p{0.40\textwidth}p{0.27\textwidth}@{}}
\hline
\textbf{Axiom} & \textbf{Claim} & \textbf{Discharge route} \\
\hline
\code{AX\_thrice\_\-punctured\_\-uniformization} & $\lambda$ realizes $\HH$ as the universal cover of $\C\setminus\{0,1\}$ ($\Gamma(2)$ deck group). & conformal-mapping programme (Schwarz reflection $\Rightarrow$ RMT $\Rightarrow$ the cover) \\
\code{AX\_developing\_\-map\_rigidity} & Two developing maps of that cover with the same monodromy agree up to a modular frame change (lift uniqueness). & classical; Forster, Ahlfors \\
\code{AX\_jacobi\_quartic} & $\theta_3^4=\theta_2^4+\theta_4^4$ (Jacobi quartic identity). & Jacobi triple product \\
\code{AX\_theta3\_ne\_zero} & $\theta_3\neq0$ on $\HH$. & Jacobi triple product \\
\code{AX\_picard\_lefschetz\_\-local} (with \code{localMonodromy}) & The family's local monodromy is the Picard--Lefschetz transvection. & Gauss--Manin / local-system layer over the moduli space \\
\hline
\end{tabular}
\end{center}

\noindent The $SU(2)$ uniqueness theorem \code{sw\_su2\_unique} rests on exactly the two B.2
covering/lift axioms plus the physics \code{SameSWMonodromy} (and the three logical axioms)---
the concrete instance of the \S4.2 split. The three logical axioms (\code{propext},
\code{Classical.choice}, \code{Quot.sound}) underlie everything and are counted in neither list.

\section*{Acknowledgments}
This formalization and paper were developed with heavy use of Claude (Anthropic), and with
Gemini (Google) and Codex (OpenAI) for cross-checking, second opinions, and review. I thank
David Kosower and Jesse Thaler for discussions on the validation of amplitude computations for
collider physics, which informed the broader perspective on validation taken here.

\section*{References}
\small
\noindent
Seiberg, Witten, \emph{Nucl.\ Phys.}\ B426 (1994) [hep-th/9407087]; B431 (1994)
[hep-th/9408099].
Klemm, Lerche, Yankielowicz, Theisen, \emph{Phys.\ Lett.}\ B344 (1995).
Argyres, Faraggi, \emph{Phys.\ Rev.\ Lett.}\ 74 (1995).
Argyres, Douglas, \emph{Nucl.\ Phys.}\ B448 (1995).
Lerche, hep-th/9611190.
Bilal, hep-th/9601007.
Nekrasov, \emph{Adv.\ Theor.\ Math.\ Phys.}\ 7 (2003); Nekrasov, Okounkov (2006).
Mathlib (Lean 4).
Paulson, Isabelle/HOL \code{Riemann\_Mapping}.
M.\ R.\ Douglas, \emph{formalization-assurance} (\code{math-commons/\discretionary{}{}{}formalization-assurance}).
Lean FRO, Tau Ceti.

\end{document}
